\documentclass[11pt]{article}
\usepackage{listings}
\usepackage{booktabs}
\usepackage{amsmath}
\usepackage{multirow}
\usepackage{minted}

% Change "review" to "final" to generate tbhe final (sometimes called camera-ready) version.
% Change to "preprint" to generate a non-anonymous version with page numbers.
\usepackage[preprint]{acl}

% Standard package includes
\usepackage{times}
\usepackage{latexsym}

% For proper rendering and hyphenation of words containing Latin characters (including in bib files)
\usepackage[T1]{fontenc}
% For Vietnamese characters
% \usepackage[T5]{fontenc}
% See https://www.latex-project.org/help/documentation/encguide.pdf for other character sets

% This assumes your files are encoded as UTF8
\usepackage[utf8]{inputenc}

% This is not strictly necessary, and may be commented out,
% but it will improve the layout of the manuscript,
% and will typically save some space.
\usepackage{microtype}

% This is also not strictly necessary, and may be commented out.
% However, it will improve the aesthetics of text in
% the typewriter font.
%\usepackage{inconsolata}

%Including images in your LaTeX document requires adding
%additional package(s)
\usepackage{graphicx}

% If the title and author information does not fit in the area allocated, uncomment the following
%
%\setlength\titlebox{<dim>}
%
% and set <dim> to something 5cm or larger.

\title{Diseño de un Benchmark y Framework de Evaluación para Robustez Adversarial en Tareas de Razonamiento Verbal en Español}

% Author information can be set in various styles:
% For several authors from the same institution:
% \author{Author 1 \and ... \and Author n \\
%         Address line \\ ... \\ Address line}
% if the names do not fit well on one line use
%         Author 1 \\ {\bf Author 2} \\ ... \\ {\bf Author n} \\
% For authors from different institutions:
% \author{Author 1 \\ Address line \\  ... \\ Address line
%         \And  ... \And
%         Author n \\ Address line \\ ... \\ Address line}
% To start a separate ``row'' of authors use \AND, as in
% \author{Author 1 \\ Address line \\  ... \\ Address line
%         \AND
%         Author 2 \\ Address line \\ ... \\ Address line \And
%         Author 3 \\ Address line \\ ... \\ Address line}

\author{
    Rodrigo Jose Alva Sáenz \\ \texttt{\small rodrigo.alva@unmsm.edu.pe} \And
    Alwin Edu Dávila Raffo \\ \texttt{\small alwin.davila@unmsm.edu.pe} \AND
    Damaris Marian Del Carpio Martinez \\ \texttt{\small damaris.delcarpio@unmsm.edu.pe} \And
    Jesus Stevan Diaz Ingol \\ \texttt{\small jesus.diaz10@unmsm.edu.pe} \AND
    Paolo Luis Flores Cóngora \\ \texttt{\small paolo.flores2@unmsm.edu.pe}
}

%\author{
%  \textbf{First Author\textsuperscript{1}},
%  \textbf{Second Author\textsuperscript{1,2}},
%  \textbf{Third T. Author\textsuperscript{1}},
%  \textbf{Fourth Author\textsuperscript{1}},
%\\
%  \textbf{Fifth Author\textsuperscript{1,2}},
%  \textbf{Sixth Author\textsuperscript{1}},
%  \textbf{Seventh Author\textsuperscript{1}},
%  \textbf{Eighth Author \textsuperscript{1,2,3,4}},
%\\
%  \textbf{Ninth Author\textsuperscript{1}},
%  \textbf{Tenth Author\textsuperscript{1}},
%  \textbf{Eleventh E. Author\textsuperscript{1,2,3,4,5}},
%  \textbf{Twelfth Author\textsuperscript{1}},
%\\
%  \textbf{Thirteenth Author\textsuperscript{3}},
%  \textbf{Fourteenth F. Author\textsuperscript{2,4}},
%  \textbf{Fifteenth Author\textsuperscript{1}},
%  \textbf{Sixteenth Author\textsuperscript{1}},
%\\
%  \textbf{Seventeenth S. Author\textsuperscript{4,5}},
%  \textbf{Eighteenth Author\textsuperscript{3,4}},
%  \textbf{Nineteenth N. Author\textsuperscript{2,5}},
%  \textbf{Twentieth Author\textsuperscript{1}}
%\\
%\\
%  \textsuperscript{1}Affiliation 1,
%  \textsuperscript{2}Affiliation 2,
%  \textsuperscript{3}Affiliation 3,
%  \textsuperscript{4}Affiliation 4,
%  \textsuperscript{5}Affiliation 5
%\\
%  \small{
%    \textbf{Correspondence:} \href{mailto:email@domain}{email@domain}
%  }
%}

\begin{document}

\maketitle

\begin{abstract}
La robustez adversarial de los modelos de lenguaje de gran escala ante
perturbaciones multilingües permanece como un área poco explorada en
español. Presentamos un benchmark de 2{,}215 ítems de razonamiento
verbal y un framework modular de evaluación para medir la robustez de
LLMs frente a cinco tipos de perturbaciones adversariales en siete
tareas. El sistema integra tres proyectos complementarios:
perturbaciones léxicas (sinónimos, pares mínimos, eliminación de
atajos), parafrásticas y \textit{cross-lingual} mediante traducción a
cinco idiomas (árabe, chino, francés, japonés, ruso). Evaluamos siete
modelos de distintas familias y tamaños, reportando siete métricas de
robustez: accuracy, accuracy drop ($\Delta$), flip rate, consistencia,
transferencia positiva, transferencia negativa y rank consistency
(correlación de Spearman sobre logprobs). Los ataques cross-lingual
producen mayor degradación que los léxicos en todos los modelos
funcionales, con el árabe causando el mayor daño ($\Delta = 7.99\%$
en qwen2.5-coder). El modelo con rendimiento cercano al azar
(llama3.2-uncensored, accuracy basal 22.66\%) \textit{mejora} bajo
todos los ataques cross-lingual ($\Delta$ negativo en los cinco
idiomas), sugiriendo que la perturbación no degrada lo que ya es
esencialmente selección aleatoria. Los valores de rank consistency
calculados sobre muestras con logprobs suficientes ($n \geq 30$) se
ubican cerca de cero, sin evidencia de inversión sistemática del
ranking interno de opciones.
\end{abstract}

% ─────────────────────────────────────────────────────
\section{Introducción}

Los modelos de lenguaje de gran escala han alcanzado un rendimiento
sobresaliente en numerosas tareas de procesamiento de lenguaje natural,
incluyendo comprensión lectora, razonamiento lógico y generación de
texto \citep{brown2020language, wei2022emergent}. Sin embargo, gran
parte de la investigación sobre su evaluación se ha centrado
desproporcionadamente en el inglés, dejando subrepresentadas lenguas
con cientos de millones de hablantes \citep{multilingual2023}.
 
Un problema fundamental es la falta de estandarización: distintos
estudios utilizan datasets heterogéneos, métricas incomparables y
configuraciones experimentales inconsistentes, lo que dificulta la
comparación justa entre modelos \citep{liang2022holistic}. A esto se
suma la fragilidad adversarial: pequeñas perturbaciones en la entrada
pueden provocar fallos inesperados incluso en modelos de alto
rendimiento \citep{wang2023promptrobust}.
 
Este trabajo es el Proyecto 3 de una serie de tres proyectos
complementarios. Los Proyectos 1 y 2 construyeron datasets adversariales
mediante perturbaciones léxicas y parafrásticas. El presente integra
esos resultados bajo un framework unificado, añade perturbaciones
cross-lingual a cinco idiomas, e implementa métricas de robustez
inspiradas en \citet{languagespecific2025}, quienes estudian cómo los
procesos latentes específicos de cada lengua limitan la transferencia
cross-lingual en LLMs.
 
Las contribuciones principales son: (1) un dataset de 2{,}215 ítems
de razonamiento verbal en español bajo formato unificado y validado;
(2) un framework modular extensible que soporta
ataque cross-lingual; (3) siete métricas de robustez, incluyendo rank consistency
sobre logprobs; y (4) un análisis de patrones de fallo sobre siete
modelos bajo condiciones adversariales controladas.


\section{Problemática}

La evaluación rigurosa de la robustez adversarial en LLMs enfrenta
tres obstáculos fundamentales.
 
\paragraph{Falta de estandarización.}
Los estudios existentes utilizan datasets con formatos propietarios,
métricas inconsistentes y protocolos no reproducibles
\citep{liang2022holistic}, haciendo que los resultados publicados sean
en la práctica incomparables entre sí.
 
\paragraph{Ausencia de análisis cross-lingual en español.}
Benchmarks como PromptBench \citep{wang2023promptrobust} o AdvGLUE
\citep{wang2022adversarial} operan casi exclusivamente en inglés. La
robustez de LLMs frente a cambios de idioma en tareas de razonamiento
verbal en español, en particular para lenguas tipológicamente diversas
como el árabe, el japonés o el ruso, no ha sido estudiada de forma
sistemática.
 
\paragraph{Métricas incapaces de distinguir tipos de degradación.}
Las métricas estándar (accuracy, F1) no capturan si el modelo
\textit{mantiene} su conocimiento correcto o \textit{transfiere sus
errores} al pasar de la condición base a la adversarial. La rank
consistency \citep{languagespecific2025} permite detectar
reorganizaciones en la distribución de confianza incluso cuando la
respuesta final no cambia, revelando vulnerabilidades invisibles para
métricas de accuracy sola.
 

% ─────────────────────────────────────────────────────
\section{Experimental Setup}
 
\subsection{Dataset}

 El dataset base (\texttt{processed\_dataset.json}) contiene 2{,}215
ítems de razonamiento verbal en español (la totalidad del banco de
preguntas del proyecto), distribuidos en siete tipos de tarea que
cubren distintas dimensiones del procesamiento lingüístico
(Tabla~\ref{tab:dataset_stats}).
 
\begin{table}[h]
  \centering\small
  \begin{tabular}{lr}
    \toprule
    \textbf{Tarea} & \textbf{Ítems} \\
    \midrule
    Comprensión lectora (\texttt{reading\_comprehension})   & 500 \\
    Series verbales (\texttt{verbal\_series})               & 500 \\
    Oraciones eliminadas (\texttt{sentence\_elimination})   & 500 \\
    Ordenamiento oraciones (\texttt{sentence\_ordering})    & 319 \\
    Analogías (\texttt{analogies})                          & 209 \\
    Sinónimos/antónimos (\texttt{synonyms\_and\_antonyms})  & 174 \\
    Oraciones incompletas (\texttt{incomplete\_sentences})  &  13 \\
    \midrule
    \textbf{Total} & \textbf{2{,}215} \\
    \bottomrule
  \end{tabular}
  \caption{Distribución de ítems por tarea en el dataset base.}
  \label{tab:dataset_stats}
\end{table}
 
Todos los ítems siguen el esquema JSON con índice de respuesta base-0:

\begin{minted}[fontsize=\small,breaklines]{json}
{
  "id": 0,
  "task": "sentence_ordering",
  "question": "Cuál es el orden lógico...",
  "options": ["0) ...", "1) ...",
              "2) ...", "3) ..."],
  "answer": 3,
  "rationale": "Explicación del razonamiento"
}
\end{minted} 

El campo \texttt{rationale} es opcional y no se expone al modelo. El
framework valida automáticamente la integridad de cada dataset y alinea
los datasets adversariales con el baseline antes de calcular métricas.


\paragraph{Familia 1: Ataques Léxicos (Proyecto 1).}
 
La intuición es que un modelo que razona semánticamente, no sobre
patrones de superficie, debe mantener su respuesta ante variaciones
léxicas que no alteran el significado.
 
\begin{itemize}
  \item \textbf{Sinónimos (\texttt{Synonym}):} Palabras clave se
  reemplazan por sinónimos de frecuencia equivalente. Evalúa si el
  modelo ancla su razonamiento a tokens específicos.
 
  \item \textbf{Pares mínimos (\texttt{MinimalPairs}):} Se modifica
  una única palabra crítica (negación, cuantificador, conector) que
  invierte o altera el significado. Aísla la sensibilidad a elementos
  de alta carga semántica con mínimo cambio de superficie.
 
  \item \textbf{Eliminación de atajos (\texttt{ShortcutRemoval}):}
  Se eliminan conectores lógicos explícitos («porque», «por lo tanto»,
  «primero... luego»). Evalúa si los modelos explotan estas señales
  superficiales en lugar de construir una representación coherente.
\end{itemize}
 
\paragraph{Familia 2: Ataque Parafrástico (Proyecto 2).}
 
\begin{itemize}
  \item \textbf{Paráfrasis (\texttt{Paraphrasing}):} El enunciado se
  reescribe preservando significado pero alterando forma y estructura
  sintáctica. Más agresivo que la sustitución léxica, puede cambiar
  el orden de la información y la organización oracional completa.
\end{itemize}
 
\paragraph{Familia 3: Ataque Cross-lingual (Proyecto 3).}
 
Motivados por \citet{languagespecific2025}, evaluamos si los procesos
latentes del idioma de la instrucción interfieren con el razonamiento.
El ataque \texttt{CrossLingual} traduce el campo \texttt{question}
dinámicamente a otro idioma mediante qwen3.6-plus (OpencodeGo),
manteniendo las opciones en español. Se aplica a cinco idiomas con
distancias tipológicas crecientes respecto al español:
 
\begin{itemize}
  \item \textbf{Francés (FR):} Lengua romance; distancia tipológica
  baja. Control de mínimo impacto esperado.
  \item \textbf{Ruso (RU):} Lengua eslava; escritura cirílica,
  morfología flexiva. Distancia media-alta.
  \item \textbf{Chino (ZH):} Lengua sino-tibetana; escritura
  logográfica, morfología aislante. Distancia alta.
  \item \textbf{Japonés (JA):} Escritura mixta (hiragana, katakana,
  kanji). Distancia alta.
  \item \textbf{Árabe (AR):} Lengua semítica; escritura RTL, morfología
  radicalmente diferente. Distancia alta.
\end{itemize}
 
Las traducciones se validan contra un umbral de longitud mínima del
23\% del texto original.

% ─────────────────────────────────────────────────────
\subsection{Métricas de Evaluación}
 
Se definen siete métricas. Las de robustez son dirigidas y se calculan
para cada par ataque-vs-baseline y para pares cruzados entre ataques.
 
\paragraph{Accuracy.}
$\text{Acc} = \text{correctos} / \text{total}$.
 
\paragraph{Accuracy Drop ($\Delta$).}
\begin{equation}
  \Delta = \text{Acc}_{\text{base}} - \text{Acc}_{\text{adv}}
\end{equation}
Valor positivo indica degradación; negativo indicaría mejora aparente
(posiblemente artefacto de tasa de fallo diferencial, ver \ref{sec:analisis}).
 
\paragraph{Flip Rate.}
Proporción de aciertos en la condición base que se convierten en errores
bajo perturbación.
 
\paragraph{Consistencia.}
Proporción de ítems con la misma respuesta en ambas condiciones,
correcta o no.
 
\paragraph{Transferencia Positiva (PT).}
Fracción de aciertos base que se mantienen bajo perturbación.
 
\paragraph{Transferencia Negativa (NT).}
Fracción de errores base que se repiten \textit{idénticamente} bajo
perturbación (errores persistentes).
Un NT alto indica que el modelo no solo falla, sino que falla de
la misma manera sistemática bajo perturbación: señal de sesgo
estructural más que sensibilidad al idioma.
 
\paragraph{Rank Consistency ($\rho$).}
Correlación de Spearman entre vectores de ranking de opciones
construidos a partir de logprobs \citep{languagespecific2025}.
Captura reordenamientos en la distribución de confianza incluso cuando
la respuesta final no cambia. Requiere que el proveedor devuelva
logprobs por token y un mínimo de 30 muestras con logprobs
simultáneamente en baseline y ataque; cuando esta condición no se
cumple, se reporta como «—». En esta evaluación, logprobs se habilitó
para nidum-gemma-3, qwen2.5-coder, dolphin3, llama3.2-uncensored y
qwen3.5-plus, con cobertura variable debido al muestreo parcial.
gemma2:27b y deepseek-v4-flash no disponen de logprobs.
 

% ─────────────────────────────────────────────────────
\section{Resultados Obtenidos}


\subsection{Resultados Generales: Accuracy Base}
 
La Tabla~\ref{tab:baseline} presenta la accuracy baseline de los siete
modelos evaluados sobre el dataset completo (2{,}215 ítems).
 
\begin{table*}[t]
  \centering\small
  \begin{tabular}{lcrrrrrrrr}
    \toprule
    \textbf{Modelo} & \textbf{n} &
    \textbf{Acc} &
    \textbf{read} & \textbf{verb} & \textbf{syn} &
    \textbf{ord} & \textbf{eli} & \textbf{ana} \\
    \midrule
    deepseek-v4-flash     & 1734 & \textbf{81.55\%} & 82.79\% & 85.48\% & — & 88.03\% & 72.12\% & 81.60\% \\
    qwen3.5-plus (pool:2) & 886  & \textbf{74.27\%} & 83.12\% & — & — & 85.89\% & 65.31\% & — \\
    gemma2:27b            & 2215 & \textbf{53.14\%} & 74.60\% & 62.80\% & 58.05\% & 43.89\% & 30.40\% & 42.58\% \\
    nidum-gemma-3-27b     & 2215 & \textbf{42.62\%} & 59.40\% & 51.80\% & 50.57\% & 32.92\% & 25.80\% & 29.19\% \\
    qwen2.5-coder-7b      & 2215 & \textbf{39.64\%} & 60.40\% & 43.20\% & 42.53\% & 27.90\% & 25.40\% & 31.10\% \\
    dolphin3              & 2215 & \textbf{31.20\%} & 48.80\% & 31.60\% & 34.48\% & 23.82\% & 21.00\% & 22.01\% \\
    llama3.2-uncensored   & 2215 & \textbf{22.66\%} & 25.80\% & 26.60\% & 20.11\% & 19.75\% & 18.40\% & 23.44\% \\
    \bottomrule
  \end{tabular}
  \caption{Accuracy base por modelo y tarea. deepseek y qwen3.5-plus tienen cobertura
  parcial debido a nuestros límites de tiempo y compute. read = reading\_comprehension,
  verb = verbal\_series, syn = synonyms\_and\_antonyms,
  ord = sentence\_ordering, eli = sentence\_elimination, ana = analogies.}
  \label{tab:baseline}
\end{table*}
 
Los modelos de nube (deepseek-v4-flash, qwen3.5-plus) lideran el
rendimiento base (74--82\%), pero solo disponen de evaluación baseline
parcial debido a límites de tiempo y compute. De los modelos con
evaluación adversarial completa, gemma2:27b es el
más capaz (53.14\%) con cero fallidos de parseo. llama3.2-uncensored
obtiene 22.66\% con un solo fallo de parseo: su
bajo rendimiento es de razonamiento, no de formato, a diferencia de
lo observado en ejecuciones previas donde fallos de formato inflaban
artificialmente la tasa de errores.
 
\subsection{Resultados por Proyecto}
 
\paragraph{Proyectos 1 y 2: Ataques Léxicos y Parafrástico.}
 
La Tabla~\ref{tab:p1p2} presenta los resultados de los ataques léxicos
y parafrástico (accuracy base independiente = 61.82\%).
 
\begin{table}[h]
  \centering\small
  \begin{tabular}{lccr}
    \toprule
    \textbf{Ataque} & \textbf{Int.} & \textbf{Acc.} & \textbf{$\Delta$ \quad FR} \\
    \midrule
    Baseline              & —   & 61.82\% & — \\
    \midrule
    \multirow{3}{*}{MinimalPairs}
      & 0.1 & 59.60\% & +2.24\% \enspace 11.10\% \\
      & 0.5 & 58.17\% & +3.65\% \enspace 13.45\% \\
      & 1.0 & 58.17\% & +3.66\% \enspace 13.33\% \\
    \midrule
    Paraphrasing          & 0.5 & 58.79\% & +3.03\% \enspace 16.10\% \\
    \midrule
    \multirow{3}{*}{ShortcutRemoval}
      & 0.1 & 61.61\% & +0.22\% \enspace \phantom{0}6.85\% \\
      & 0.5 & 61.50\% & +0.32\% \enspace \phantom{0}7.61\% \\
      & 1.0 & 61.62\% & +0.22\% \enspace \phantom{0}7.49\% \\
    \midrule
    Synonym               & 0.5 & 56.80\% & +2.66\% \enspace 15.95\% \\
    \bottomrule
  \end{tabular}
  \caption{Resultados de ataques léxicos y parafrástico (P1/P2).
  FR = flip rate. Evaluación independiente sobre subconjunto del dataset.}
  \label{tab:p1p2}
\end{table}
 
\texttt{ShortcutRemoval} produce el impacto más bajo de toda la
evaluación ($\Delta \leq 0.32\%$ a cualquier intensidad), indicando
que los modelos no dependen de los conectores lógicos explícitos para
resolver las tareas. \texttt{Paraphrasing} presenta el flip rate más
alto (16.10\%) con $\Delta$ moderado (3.03\%): el modelo cambia sus
respuestas frecuentemente, pero no siempre en la dirección equivocada.
\texttt{MinimalPairs} satura a partir de intensidad 0.5: los resultados
en 0.5 y 1.0 son prácticamente idénticos, lo que sugiere que incluso
una sola modificación crítica agota el efecto adversarial de este ataque.
 
\paragraph{Proyecto 3: Ataques Cross-lingual.}
 
La Tabla~\ref{tab:cl_robustness} presenta las métricas de robustez
cross-lingual para los cinco modelos con evaluación adversarial.
 
\begin{table*}[t]
  \centering\small
  \begin{tabular}{llccccccc}
    \toprule
    \textbf{Modelo} & \textbf{Idioma} & \textbf{Acc.} &
    \textbf{$\Delta$} & \textbf{FR} & \textbf{Cons.} &
    \textbf{PT} & \textbf{NT} & \textbf{$\rho$} \\
    \midrule
    \multirow{5}{*}{gemma2:27b}
      & Árabe   & 47.67\% & +5.46\% & 32.63\% & 53.77\% & 67.37\% & 38.34\% & — \\
      & Ruso    & 48.80\% & +4.33\% & 27.87\% & 58.10\% & 72.13\% & 42.20\% & — \\
      & Japonés & 48.98\% & +4.15\% & 30.25\% & 54.18\% & 69.75\% & 36.51\% & — \\
      & Chino   & 50.20\% & +2.93\% & 27.61\% & 56.57\% & 72.39\% & 38.63\% & — \\
      & Francés & 51.69\% & +1.44\% & 23.70\% & 60.99\% & 76.30\% & 43.64\% & — \\
    \midrule
    \multirow{5}{*}{nidum-gemma-3}
      & Árabe   & 38.10\% & +4.51\% & 43.54\% & 42.26\% & 56.46\% & 31.71\% & — \\
      & Japonés & 40.18\% & +2.44\% & 42.27\% & 40.72\% & 57.73\% & 28.09\% & 0.13 \\
      & Ruso    & 41.13\% & +1.49\% & 39.51\% & 43.88\% & 60.49\% & 31.55\% & 0.23 \\
      & Francés & 42.21\% & +0.41\% & 37.08\% & 45.91\% & 62.92\% & 33.28\% & — \\
      & Chino   & 42.30\% & +0.32\% & 38.77\% & 43.52\% & 61.23\% & 30.37\% & — \\
    \midrule
    \multirow{5}{*}{qwen2.5-coder}
      & Árabe   & 31.65\% & +7.99\% & 49.43\% & 44.06\% & 50.57\% & 39.79\% & — \\
      & Japonés & 36.70\% & +2.93\% & 37.36\% & 52.10\% & 62.64\% & 45.18\% & — \\
      & Ruso    & 37.65\% & +1.99\% & 34.40\% & 56.48\% & 65.60\% & 50.49\% & 0.12 \\
      & Francés & 39.59\% & +0.05\% & 32.00\% & 53.00\% & 68.00\% & 43.16\% & — \\
      & Chino   & 41.35\% & $-$1.72\% & 33.03\% & 51.83\% & 66.97\% & 41.88\% & — \\
    \midrule
    \multirow{5}{*}{dolphin3}
      & Árabe   & 26.82\% & +4.38\% & 57.89\% & 34.72\% & 42.11\% & 31.36\% & 0.22 \\
      & Japonés & 27.58\% & +3.61\% & 54.27\% & 35.85\% & 45.73\% & 31.36\% & 0.14 \\
      & Ruso    & 28.08\% & +3.12\% & 52.82\% & 36.21\% & 47.18\% & 31.23\% & 0.04 \\
      & Francés & 30.02\% & +1.17\% & 46.89\% & 38.92\% & 53.11\% & 32.48\% & — \\
      & Chino   & 30.25\% & +0.95\% & 50.51\% & 35.58\% & 49.49\% & 29.27\% & — \\
    \midrule
    \multirow{5}{*}{llama3.2-uncs.}
      & Árabe   & 23.57\% & $-$0.90\% & 69.72\% & 26.28\% & 30.28\% & 25.10\% & $-$0.06 \\
      & Francés & 25.06\% & $-$2.39\% & 64.94\% & 28.44\% & 35.06\% & 26.50\% & $-$0.01 \\
      & Ruso    & 25.10\% & $-$2.44\% & 70.92\% & 25.10\% & 29.08\% & 23.93\% & $-$0.00 \\
      & Chino   & 25.24\% & $-$2.57\% & 68.13\% & 25.19\% & 31.87\% & 23.23\% & 0.11 \\
      & Japonés & 26.68\% & $-$4.02\% & 65.94\% & 26.68\% & 34.06\% & 24.52\% & $-$0.06 \\
    \bottomrule
  \end{tabular}
  \caption{Métricas de robustez cross-lingual por modelo e idioma, ordenadas
  por $\Delta$ descendente dentro de cada modelo. $\rho$ = rank consistency
  (Spearman sobre logprobs, $n \geq 30$); — = logprobs no disponibles
  o muestra insuficiente. PT = transferencia
  positiva; NT = transferencia negativa.}
  \label{tab:cl_robustness}
\end{table*}
 
La Tabla~\ref{tab:gemma_pertask} desglosa los resultados de gemma2:27b
por tarea y los cinco idiomas de perturbación.
 
\begin{table*}[t]
  \centering\small
  \begin{tabular}{lrrrrrrrrrr}
    \toprule
    & \textbf{Base} &
    \multicolumn{2}{c}{\textbf{Árabe}} &
    \multicolumn{2}{c}{\textbf{Ruso}} &
    \multicolumn{2}{c}{\textbf{Japonés}} &
    \multicolumn{2}{c}{\textbf{Chino}} &
    \multicolumn{1}{c}{\textbf{Fran.}} \\
    \textbf{Tarea} & Acc & Acc & $\Delta$ & Acc & $\Delta$ & Acc & $\Delta$ & Acc & $\Delta$ & $\Delta$ \\
    \midrule
    reading\_comp.   & 74.60\% & 67.60\% & +7.00\% & 67.80\% & +6.80\% & 68.80\% & +5.80\% & 68.00\% & +6.60\% & +3.80\% \\
    verbal\_series   & 62.80\% & 54.60\% & +8.20\% & 59.40\% & +3.40\% & 59.40\% & +3.40\% & 62.40\% & +0.40\% & +1.20\% \\
    syn./ant.        & 58.05\% & 54.60\% & +3.45\% & 55.75\% & +2.30\% & 59.20\% & $-$1.15\% & 63.79\% & $-$5.75\% & $-$3.45\% \\
    inc.\_sentences  & 61.54\% & 38.46\% & +23.08\% & 61.54\% & \phantom{+}0.00\% & 61.54\% & \phantom{+}0.00\% & 38.46\% & +23.08\% & +15.38\% \\
    analogies        & 42.58\% & 33.97\% & +8.61\% & 34.93\% & +7.66\% & 30.14\% & +12.44\% & 34.45\% & +8.13\% & +0.96\% \\
    sentence\_ord.   & 43.89\% & 47.34\% & $-$3.45\% & 44.51\% & $-$0.63\% & 46.39\% & $-$2.51\% & 45.45\% & $-$1.57\% & $-$7.84\% \\
    sentence\_eli.   & 30.40\% & 24.60\% & +5.80\% & 25.00\% & +5.40\% & 24.40\% & +6.00\% & 25.40\% & +5.00\% & +6.80\% \\
    \midrule
    \textbf{Total}   & \textbf{53.14\%} & \textbf{47.67\%} & \textbf{+5.46\%} & \textbf{48.80\%} & \textbf{+4.33\%} & \textbf{48.98\%} & \textbf{+4.15\%} & \textbf{50.20\%} & \textbf{+2.93\%} & \textbf{+1.44\%} \\
    \bottomrule
  \end{tabular}
  \caption{Accuracy y accuracy drop por tarea, gemma2:27b (5 idiomas).
  $\Delta$ negativo = mejora bajo ataque. Para francés, solo se muestra
  $\Delta$ por espacio; accuracy base $\times$ (1$-\Delta$) para
  recuperar el valor.}
  \label{tab:gemma_pertask}
\end{table*}
 
% ─────────────────────────────────────────────────────
\subsection{Análisis de Resultados}
\label{sec:analisis}

\paragraph{Patrón 1: Árabe produce el mayor daño; francés el menor.}
En los cuatro modelos con mayor capacidad (gemma2:27b, nidum-gemma-3,
qwen2.5-coder, dolphin3), el árabe es el idioma de ataque más dañino
(gemma2:27b: $\Delta = +5.46\%$; nidum-gemma-3: $\Delta = +4.51\%$;
qwen2.5-coder: $\Delta = +7.99\%$; dolphin3: $\Delta = +4.38\%$) y el
francés o chino el menos dañino (gemma2:27b: $\Delta = +1.44\%$;
nidum-gemma-3: $\Delta = +0.32\%$ chino). El ordenamiento en gemma2:27b
es: AR $>$ RU $>$ JA $>$ ZH $>$ FR, no estrictamente coincidente con la
distancia tipológica al español. El árabe supera al chino pese a que el
chino es tipológicamente más distante, posiblemente porque su escritura
RTL introduce ambigüedades adicionales en la tokenización o porque los
corpus de preentrenamiento en árabe tienen menor representación que los
de chino. llama3.2-uncensored constituye una excepción a este patrón
(ver Patrón 3).
 
\paragraph{Patrón 2: Los ataques cross-lingual superan a los léxicos.}
El $\Delta$ máximo de los ataques léxicos y parafrástico (P1/P2) es
3.66\% (\texttt{MinimalPairs} intensidad 1.0). El ataque árabe produce
$+4.38\%$ a $+7.99\%$ según el modelo, es decir, entre 1.2× y 2.2× mayor
degradación. \texttt{ShortcutRemoval} representa el extremo inferior:
$\Delta \leq 0.32\%$ a cualquier intensidad, lo que indica que los
conectores lógicos explícitos contribuyen marginalmente al rendimiento
en estas tareas.
 
\paragraph{Patrón 3: llama3.2-uncensored mejora bajo todos los ataques
cross-lingual.}
llama3.2-uncensored es el único modelo donde los cinco idiomas de
perturbación producen $\Delta$ \textit{negativo}: AR ($-$0.90\%),
FR ($-$2.39\%), RU ($-$2.44\%), ZH ($-$2.57\%), JA ($-$4.02\%). Es
decir, el modelo obtiene \textit{mayor} accuracy bajo ataque que en
baseline. Esto no indica robustez, sino lo contrario: con una accuracy
basal de 22.66\%, cercana al azar (25\% para 4 opciones, 20\% para 5),
el modelo no tiene un rendimiento significativo que degradar. Las
flip rates altas (65--71\%) confirman que el modelo cambia sus
respuestas con alta frecuencia, pero aproximadamente en igual proporción
hacia aciertos y hacia errores. Este patrón es consistente con
selección casi aleatoria de opciones, donde la perturbación no
empeora ni mejora un razonamiento inexistente.
 
\paragraph{Patrón 4: Rank consistency cercana a cero.}
Los valores de $\rho$ calculados sobre muestras con logprobs
suficientes ($n \geq 30$) se ubican cerca de cero en todos los casos.
El único modelo con cobertura completa de $\rho$ en los cinco idiomas
es llama3.2-uncensored, con valores entre $-$0.06 y 0.11. Esto es
coherente con su comportamiento cercano al azar: sin un ranking
estable de opciones que preservar, la correlación entre baseline y
ataque es nula. Los valores más altos de $\rho$ se observan en
nidum-gemma-3/ruso ($\rho = 0.23$, $n = 162$) y dolphin3/árabe
($\rho = 0.22$, $n = 140$), indicando correlación positiva moderada:
el modelo mantiene parcialmente su ranking bajo perturbación. En
qwen2.5-coder/ruso, $\rho = 0.12$ ($n = 212$), cercano al cero. Los
valores de $\rho$ para varios pares modelo-idioma no pudieron
calcularse debido a cobertura insuficiente de logprobs ($n < 30$),
reportándose como «—» en la Tabla~\ref{tab:cl_robustness}.
 
\paragraph{Patrón 5: sentence\_ordering mejora consistentemente bajo
perturbación cross-lingual.}
En gemma2:27b, \texttt{sentence\_ordering} mejora bajo los cinco idiomas
de perturbación: $\Delta = -3.45\%$ (AR), $-0.63\%$ (RU), $-2.51\%$
(JA), $-1.57\%$ (ZH), $-7.84\%$ (FR). Una hipótesis es que cuando la
instrucción llega en otro idioma, el modelo reduce el procesamiento
semántico profundo y se enfoca en señales estructurales: en
\texttt{sentence\_ordering}, el orden lógico puede inferirse de
marcadores de cohesión (pronombres, tiempos verbales) que son
independientes del idioma de la instrucción. También se observan mejoras
menores en \texttt{synonyms\_and\_antonyms} bajo chino ($-5.75\%$) y
francés ($-3.45\%$).
 
\paragraph{Patrón 6: Alta flip rate con $\Delta$ negativo en
llama3.2-uncensored.}
Como se observa en el Patrón 3, llama3.2-uncensored presenta flip rates
altas (65--71\%) bajo todos los ataques, con $\Delta$ \textit{negativos}
en los cinco idiomas ($-$0.90\% a $-$4.02\%). La combinación de alta
flip rate con $\Delta$ negativo confirma que el modelo cambia sus
respuestas frecuentemente, pero con un ligero sesgo hacia aciertos bajo
perturbación. El $\rho \approx 0$ (entre $-$0.06 y 0.11) bajo los cinco
idiomas indica que no hay un ranking interno estable que la perturbación
pueda reordenar, consistente con selección casi aleatoria.
 
\paragraph{Patrón 7: Transferencia negativa alta como indicador de sesgo.}
En gemma2:27b, el NT oscila entre 36.5\% y 43.6\% según el idioma,
indicando que entre 37\% y 44\% de los errores en baseline se repiten
\textit{idénticamente} bajo perturbación. En qwen2.5-coder, el NT
bajo ruso alcanza 50.5\%: la mitad de todos los errores del modelo son
persistentes e independientes del idioma del ataque. Esto sugiere
sesgos estructurales preexistentes, como preferencia sistemática por
ciertas posiciones de opción, que ninguna perturbación logra alterar.
 
% ─────────────────────────────────────────────────────
\subsection{Interfaces e Implementación}
 
El framework \texttt{llm\_verbal\_framework} se organiza en cuatro
capas.
 
\paragraph{Capa de modelos (Providers).}
Tres backends bajo una interfaz común
(\texttt{complete(messages, response\_format) $\to$ content, tokens,
logprobs}): Ollama (modelos locales), OpenRouter (modelos propietarios
con logprobs) y OpencodeGo (dual backend: AsyncOpenAI para kimi-k2.6,
deepseek-v4-flash, qwen3.5-plus; AsyncAnthropic para modelos Anthropic).
 
\paragraph{Capa de ataques.}
\texttt{CrossLingual} genera perturbaciones dinámicamente con validación
de longitud. Los demás ataques cargan datasets pre-computados:
 
\begin{minted}[fontsize=\small,breaklines]{python}
benchmark = Benchmark(
    baseline=Path("processed_dataset.json"),
    attacks=[
        attacks.CrossLingual(
            language=CrossLingualLanguage.ARABIC),
        attacks.CrossLingual(
            language=CrossLingualLanguage.RUSSIAN),
        attacks.CrossLingual(
            language=CrossLingualLanguage.JAPANESE),
        attacks.CrossLingual(
            language=CrossLingualLanguage.CHINESE),
        attacks.CrossLingual(
            language=CrossLingualLanguage.FRENCH),
        attacks.Synonym("synonym_1"),
        attacks.Paraphrasing("paraphrasing_1"),
        attacks.MinimalPairs("negation_flip"),
        attacks.ShortcutRemoval("no_shortcuts"),
    ],
    models=[...],
    concurrency=4, base_dir="benchmark",
)
result = benchmark.run()
result.save("stats.json")
result.save("report.md")
\end{minted}
 
\paragraph{Resumabilidad y logprobs.}
El framework guarda resultados parciales tras cada lote mediante
escritura atómica (archivo temporal + rename), permitiendo reanudar
evaluaciones interrumpidas sin re-procesar trabajo completado. Con
\texttt{logprobs=True} y \texttt{top\_logprobs=5}, el framework extrae
logprobs por opción escaneando tokens de dígito en el output del modelo
y calcula $\rho$ sobre los vectores de ranking resultantes.
 
\paragraph{Capa de reporte.}
Exportación en JSON (agregados y por muestra) y Markdown con tablas
científicas. La evaluación total duró aproximadamente 20 horas
combinando dos ejecuciones del benchmark sobre hardware local (Ollama)
y APIs en la nube.

 
% ─────────────────────────────────────────────────────
\section{Conclusiones}
Este trabajo presentó un benchmark de 2{,}215 ítems de razonamiento
verbal en español y un framework de evaluación modular para medir la
robustez adversarial de LLMs bajo cinco tipos de ataque y cinco idiomas
de perturbación. Las contribuciones y hallazgos principales son:
 
\begin{enumerate}
  \item Un dataset unificado en formato JSON con siete tipos de tarea,
  con validación automática de integridad estructural.
 
  \item Un framework reproducible con tres backends de modelo,
  concurrencia configurable y exportación en múltiples formatos.
 
  \item Siete métricas de robustez incluyendo rank consistency sobre
  logprobs, que requiere un mínimo de 30 muestras con logprobs
  simultáneos para producir un valor fiable.
 
  \item La jerarquía de impacto adversarial: ataques cross-lingual
  dominan sobre léxicos ($\Delta_{\text{AR}} \leq 7.99\%$ vs
  $\Delta_{\text{MinPairs}} \leq 3.66\%$); dentro de los cross-lingual,
  árabe $>$ ruso $\approx$ japonés $>$ chino $>$ francés en los modelos
  con mayor capacidad.
  \texttt{ShortcutRemoval} es el ataque menos efectivo ($\Delta \leq
  0.32\%$), lo que indica que los modelos no explotan estos atajos
  superficiales de manera significativa.
 
  \item Los valores de rank consistency calculados sobre muestras
  suficientes ($n \geq 30$) se ubican cerca de cero, sin evidencia de
  inversión sistemática del ranking interno de opciones. La cobertura
  parcial de logprobs limita el análisis a un subconjunto de pares
  modelo-idioma.
 
  \item llama3.2-uncensored, con accuracy basal cercana al azar
  (22.66\%), \textit{mejora} bajo todos los ataques cross-lingual
  ($\Delta$ negativo en los cinco idiomas). Esto no indica robustez
  sino ausencia de razonamiento: la perturbación no puede degradar lo
  que es esencialmente selección aleatoria.
 
  \item \texttt{sentence\_ordering} mejora consistentemente bajo
  todos los idiomas de perturbación en gemma2:27b, sugiriendo que
  la instrucción en otro idioma favorece un procesamiento más
  estructural que beneficia a esta tarea en particular.
\end{enumerate}
 
Como trabajo futuro, se propone completar la evaluación de los modelos
de nube con ataques adversariales (deepseek-v4-flash, qwen3.5-plus),
extender la cobertura de logprobs a todos los modelos y muestras para
habilitar un análisis más robusto de rank consistency, y aplicar
análisis de representaciones internas (CKA, activation steering) para
contrastar con los hallazgos de \citet{languagespecific2025}.
% ─────────────────────────────────────────────────────
\section*{Agradecimientos}
 
Los autores agradecen al PhD. Marco Antonio Sobrevilla Cabezudo, docente del curso de Software Inteligente de la Universidad Nacional Mayor de San Marcos por la orientación brindada durante el desarrollo de este trabajo.
 
% ─────────────────────────────────────────────────────
\bibliography{custom}

% ─────────────────────────────────────────────────────
\appendix
 
\section{Ejemplos de Perturbaciones por Tipo de Ataque}
\label{sec:appendix_examples}

Los ejemplos se resumen en la Tabla ~\ref{tab:appendix_examples}.
 
\begin{table*}[h]
  \centering\small
  \begin{tabular}{p{2.2cm}p{5.3cm}p{5.3cm}}
    \toprule
    \textbf{Ataque} & \textbf{Enunciado original (ES)} & \textbf{Enunciado perturbado} \\
    \midrule
    Base
      & ¿Cuál es sinónimo de «perspicaz»? & — \\
    Synonym
      & ¿Cuál es sinónimo de «sagaz»?
      & Sustitución léxica; significado preservado. \\
    MinimalPairs
      & ¿Cuál \textbf{no} es sinónimo de «perspicaz»?
      & Un cambio mínimo invierte el significado. \\
    ShortcutRemoval
      & ¿Cuál corresponde a «perspicaz»?
      & Eliminación del atajo «sinónimo de». \\
    Paraphrasing
      & ¿Qué palabra tiene el mismo significado que «perspicaz»?
      & Reescritura sintáctica completa. \\
    CL/FR (francés)
      & Lequel est synonyme de «perspicaz»?
      & Traducción al francés; opciones en ES. \\
    CL/ZH (chino)
      & 以下哪个词是"perspicaz"的同义词？
      & Traducción al chino; opciones en ES. \\
    CL/AR (árabe)
      & أيّ من المصطلحات التالية هو مرادف لكلمة «perspicaz»؟
      & Traducción al árabe (RTL); opciones en ES. \\
    CL/JA (japonés)
      & 「perspicaz」の同義語はどれですか？
      & Traducción al japonés; opciones en ES. \\
    CL/RU (ruso)
      & Какое из следующих слов является синонимом «perspicaz»?
      & Traducción al ruso (cirílico); opciones en ES. \\
    \bottomrule
  \end{tabular}
  \caption{Ejemplos de los nueve tipos de perturbación sobre un ítem de
  sinónimos/antónimos. Las opciones permanecen en español en todos los casos.}
  \label{tab:appendix_examples}
\end{table*}
 
\section{Resumen de Modelos y Cobertura}
\label{sec:appendix_models}

Los modelos evaluados se resumen en la Tabla~\ref{tab:models_summary}.
 
\begin{table*}[h]
  \centering\small
  \begin{tabular}{lcccc}
    \toprule
    \textbf{Modelo} & \textbf{Prov.} & \textbf{Ataques CL} & \textbf{logprobs} \\
    \midrule
    deepseek-v4-flash     & OCG & --- & No \\
    qwen3.5-plus (pool:2) & mix & parcial & Parcial \\
    gemma2:27b            & Oll & AR/RU/JA/ZH/FR & No \\
    nidum-gemma-3-27b     & Oll & AR/RU/JA/ZH/FR & Parcial \\
    qwen2.5-coder-7b      & Oll & AR/RU/JA/ZH/FR & Parcial \\
    dolphin3              & Oll & AR/RU/JA/ZH/FR & Parcial \\
    llama3.2-uncensored   & Oll & AR/RU/JA/ZH/FR & Parcial \\
    \bottomrule
  \end{tabular}
  \caption{Modelos evaluados, ataques ejecutados y disponibilidad
  de logprobs para rank consistency. ``Parcial'' indica que los logprobs
  se recolectaron solo para un subconjunto de muestras, por lo que
  algunos pares modelo-idioma no alcanzan el mínimo de $n = 30$ para
  $\rho$. qwen3.5-plus dispone de evaluación CL parcial
  ($n = 308$--$476$ por idioma) debido a límites de compute.
  OCG = OpencodeGo, Oll = Ollama, mix = pool Ollama+OpencodeGo.}
  \label{tab:models_summary}
\end{table*}
 
\end{document}